% Tabellenkarten für den Werbe-Teil von InstaHub (Spalten gemäss den InstaHub-Migrationen,
% https://github.com/wi-wissen/InstaHub, database/migrations/create/ads und .../analytics).
% Jede Tabelle ist eine einspaltige Matrix: Zeile 1 = Tabellenname, danach die Kolonnen.
% Die Pfeile zeigen, auf welche id die Kolonnen user_id und photo_id verweisen.
\begin{tikzpicture}[
    karte/.style={
        matrix of nodes, draw, thick, inner sep=0pt, row sep=0pt,
        nodes={anchor=west, font=\scriptsize\ttfamily, inner sep=1.5pt, text height=1.6ex, text depth=.4ex, minimum width=26mm},
        row 1/.style={nodes={font=\scriptsize\bfseries\ttfamily, fill=RoyalBlue!20}},
      },
    verweis/.style={->, >=stealth, thick, RoyalBlue},
  ]
  \matrix (ads) [karte] {
    ads \\ id \\ priority \\ name \\ type \\ url \\ img \\ query \\ created\_at \\ updated\_at \\
  };
  \matrix (analytics) [karte, right=8mm of ads.north east, anchor=north west] {
    analytics \\ id \\ ip \\ device \\ brand\_family \\ brand\_model \\ browser\_family \\ browser\_version \\
    platform\_family \\ platform\_version \\ user\_id \\ photo\_id \\ created\_at \\ updated\_at \\
  };
  \matrix (users) [karte, right=14mm of analytics.north east, anchor=north west] {
    users \\ id \\ username \\ \ldots \\
  };
  \matrix (photos) [karte, right=14mm of analytics.south east, anchor=south west] {
    photos \\ id \\ user\_id \\ description \\ \ldots \\
  };
  % Verweise analytics.user_id -> users.id und analytics.photo_id -> photos.id
  \draw[verweis] (analytics-11-1.east) -- ++(6mm,0) |- (users-2-1.west);
  \draw[verweis] (analytics-12-1.east) -- ++(9mm,0) |- (photos-2-1.west);
\end{tikzpicture}
